
\section{[附录一] 常见文言多义词义项举要}

\textbf{1. 去}

①\ 离开：登斯楼也，则有去国怀乡，忧谗畏讥，满目萧然，感极而悲者矣。（《岳阳楼记》）
②\ 距离：当是时，项羽军在鸿门，沛公军在霸上，相去四十里。（《鸿门宴》）

③\ 去掉，除掉：去死肌，杀三虫。（《捕蛇者说》）

④\ 消失：四五日外，色香味尽去矣。（《荔枝图序》）

\textbf{2. 从}

①\ 跟随：赵王遂行，相如从。（《廉颇蔺相如列传》）

②\ 顺从：小惠未遍，民弗从也。（《曹刿论战》）

③\ 依从，听从。公子从其计，请如姬。（《信陵君窃符救赵》）

④\ 随行，侍从，读zòng：忠之属也，可以一战。战则请从。（《曹刿论战》）

⑤\ 堂房亲属，升死，其印为予群从所得，至今保藏。（《活板》）

⑥\ 由，自：从小丘西行百二十步……（《至小丘西小石潭记》）

\textbf{3. 使}

①\ 命令，派遣：秦王使使者告赵王，欲与王为好，会于西河外渑池。（《廉颇蔺相如列传》）

②\ 令，让，叫：使天下之人，不敢言而敢怒。（《阿房宫赋》）

③\ 奉使命出使，读shì：计未定，求人可使报秦者，未得。（《廉颇蔺相如列传》）

④\ 奉使命的人，使臣：会使辙交驰，北邀当国者相见……（《指南录后序》）

⑤\ 假使：使六国各爱其人，则足以拒秦。（《阿房宫赋》）

\textbf{4. 谢}

①\ 道歉：入而徐趋，至而自谢。（《触龙说赵太后》）

②\ 辞别：侯生视公子色终不变，乃谢客就车。（《信陵君窃符救赵》）

③\ 告诉：多谢后世人，戒之慎勿忘！（《孔雀东南飞》）

④\ 道谢：哙拜谢，起，立而饮之。（《鸿门宴》）

\textbf{5. 遗}

①\ 丢失：小学而大遗，吾未见其明也。（《师说》）

②\ 遗留：初患烈遗言：“我死当葬梅花岭上。”（《梅花岭记》）

③\ 遗传下来的：诚宜开张圣听，以光先帝遗德。（《出师表》）

④\ 送，读 wèi：公子姊……数遗魏王及公子书，请教于魏。（《信陵君窃符救赵》）

⑤\ 赠送，读 wèi：人贱物亦鄙，不足迎后人。留待作遗施，于今无会因。（《孔雀东南飞》）

\textbf{6. 策}

①\ 竹制的马鞭：振长策而御宇内。（《过秦论》）

②\ 鞭策、鞭打：策之不以其道，食之不能尽其材。（《马说》）

③\ 骑上鞭打：乃蹶蒙氏衣白衣，复自策其马，麾众拥蒙民马前……（《书博鸡者事》）

④\ 记在简策上：策勋十二转，赏赐百千强。（《木兰诗》）

⑤\ 策划、谋划：请为大王策之。（《左传》）

⑥\ 计策、计谋：均之二策，宁许以负秦曲。（《廉颇蔺相如列传》）

\textbf{7. 间}

①\ 间隙、夹缝：彼节者有间，而刀刃者无厚。（《庖丁解牛》）

②\ 隔离、离间：屈平正道直行……谗人间之，可谓穷矣。（《屈原列传》）

③\ 参与，置身其中；肉食者谋之，又何间焉。（《曹刿论战》）

④\ 间或，断断续续地：数月之后，时时而间进。（《邹忌讽齐王纳谏》）

⑤\ 秘密地，暗地里：又间令吴广之次所旁丛祠中，夜篝火，狐鸣呼曰……（《陈涉世家》）

⑥\ 抄近路，抄小路：从鄱山下，道芷阳间行。（《鸿门宴》）

⑦\ 中间：出师一表真名世，千载谁堪伯仲间！（《书愤》）

⑧\ 一会儿：立有间。（《扁鹊见蔡桓公》）

\textbf{8. 致}

①\ 送给：远方莫不效其珍。（《荀子·解蔽》）

②\ 得到：家贫，无从致书以观。（《送东阳马生序》）

③\ 达到：秦以区区之地，致万乘之势。（《过秦论》）

④\ 表达：一篇之中，三致意焉。（《屈原列传》）

⑤\ 招引：士以此方数千里争往归之，致食客三千人。（《信陵君窃符救赵》）

⑥\ 使至，到：假舆马者，非利足也，而致千里。（《劝学》）

⑦\ 尽，极：衡善机巧，尤致思于天文阴阳历算。（《张衡传》）

\textbf{9. 适}

①\ 到（某地）去：逝将去女，适彼乐土。（《硕鼠》）

②\ 出嫁：贫贱有此女，始适还家门。（《孔雀东南飞》）

③\ 适合、依照：处分适兄意，那得自任专。（《孔雀东南飞》）

④\ 恰好，从上观之，适与地平。（雁荡山）⑤刚才，时已过午，奴辈适至。（《游黄山记》）

\textbf{10.书}

①\ 写：乃丹书帛曰“陈胜王”。（《陈涉世家》）

②\ 文字：读其书未毕，齐军万弩俱发。（史记·吴起孙膑列传）

③\ 书信：忽报曹操遣使送书至（《赤壁之战》）

④\ 文书：昼断狱，夜理书。（《汉书·刑法志》）

⑤\ 书籍：家贫，无从致书以观。（《送东阳马生序》）

⑥\ 特指《尚书》：“《书》曰：“满招损，谦得益。”（《伶官传序》）

\textbf{11. 道}

①\ 道路：会天大雨，道不通。（《陈涉世家》）

②\ 途径、方法：此五者，知胜之道也。（《谋攻》）

③\ 道理：虽有至道，弗学不知其善也。（《学记》）

④\ 规律：臣之所好者，道也；进乎技矣。（《庖丁解牛》）

⑤\ 学说、思想：既加冠，益慕圣贤之道。（《送东阳马生序》）

⑥\ 述说：此中人语云：“不足为外人道也。”（《桃花源记》）

\textbf{12. 迁}

①\ 迁移、迁都：时北兵已迫修门外，战、守、迁皆不及施。（《指南录后序》）

②\ 变易：齐人未尝赂秦，终继五国迁灭，何哉？（《六国论》）

③\ 调动、升官：公车特征，拜郎中，再迁为太史令。（《张衡传》）

④\ 降职，贬谪；迁客骚人，多会于此。（《岳阳楼记》）

⑤\ 放逐；顷襄王怒而迁之。（《屈原列传》）

\textbf{13. 卒}

①\ 步兵；夫以疲病之卒，御狐疑之众……（《赤壁之战》）

②\ 一百人的队伍；凡用兵之法……全卒为上，破卒次之。（《谋攻》）

③\ 完毕；语未及卒，公子立变色。（《信陵君窃符救赵》）

④\ 终于；卒相与欢，为刎颈之交。（《廉颇蔺相如列传》）

⑤\ 死；初，鲁肃闻刘表卒。（《赤壁之战》）

⑥\ 突然（在这个意义上，与“猝”通，读cù）；五万兵难卒合，已选三万人……（《赤壁之战》）

\textbf{14. 举}
①\ 举起、抬起；千钧之重，人不轻举。（《盐铁论·刑德》）

②\ 升起；中江举帆，余船以次俱进。（《赤壁之战》）

③\ 举荐；举孝廉不行，连辟公府不就。（《张衡传》）

④\ 发动；今亡亦死，举大计亦死，等死，死国可乎？（《陈涉世家》）

⑤\ 攻下；戍卒叫，函谷举。（《阿房宫赋》）

⑥\ 全，尽；杀人如不能举